\section*{अध्याय \ref{ch:b1:organism-environment} --- जीव: अपने पर्यावरण से अंतःक्रिया करता एक तंत्र}

\begin{solution}{exo:b1:organism-environment:1}
पदार्थ (कोई स्तनधारी भोजन खाता है और कार्बन डाइऑक्साइड छोड़ता है; कोई
पौधा कार्बन डाइऑक्साइड लेता है और ऑक्सीजन छोड़ता है), ऊर्जा (कोई स्तनधारी
भोजन की रासायनिक ऊर्जा लेता है और ऊष्मा छोड़ता है; कोई पौधा प्रकाश लेता
है और ऊष्मा छोड़ता है), सूचना (कोई स्तनधारी किसी परभक्षी को देखता है; कोई
पौधा प्रकाश की ओर झुकता है)।
\end{solution}

\begin{solution}{exo:b1:organism-environment:2}
$S/V = 3/r$ आधा हो जाता है। पृष्ठ से ग्रहण $r^2$ की तरह बढ़ता है और
खपत $r^3$ की तरह; किसी त्रिज्या के आगे पृष्ठ आयतन को नहीं पाल सकता,
इसलिए कोशिकाएँ माइक्रोमीटर के परास में रहती हैं या अपना पृष्ठ मोड़ लेती हैं।
\end{solution}

\begin{solution}{exo:b1:organism-environment:3}
किसी नियत-बिंदु के पास आंतरिक परिवेश के किसी चर का ऋणात्मक पुनर्भरण से
बनाए रखा जाना। संवेदी: त्वचा और मस्तिष्क के तापग्राही; समाकलन केंद्र:
अधश्चेतक, जो \qty{37}{\celsius} का नियत-बिंदु रखता है; प्रभावक: कँपकँपाती
पेशियाँ, त्वचा की वाहिकाएँ, स्वेद ग्रंथियाँ।
\end{solution}

\begin{solution}{exo:b1:organism-environment:4}
छिपकली: $1.0$ से $0.25$ तक, यानी $4$ का गुणक (\qty{20}{\celsius} के लिए
दो द्विगुणन)। चूहा: $1.0$ से $3.4$ तक, यानी लगभग
$3.4$ का गुणक।
\end{solution}

\begin{solution}{exo:b1:organism-environment:5}
\qty{1}{mm} की भुजा: चाहिए $\qty{1}{nmol/s}$ ($V = \qty{1}{mm^3}$); आपूर्ति
$2\times 6 = \qty{12}{nmol/s}$: हाँ। \qty{10}{mm} भुजा: चाहिए
\qty{1000}{nmol/s}, आपूर्ति $2\times 600 = \qty{1200}{nmol/s}$: बस
किसी तरह। सीमा: $2\times 6L^2 = L^3$ से $L = \qty{12}{mm}$।
\end{solution}

\begin{solution}{exo:b1:organism-environment:6}
घोड़ा: $3.4\times 500^{0.75} = \qty{360}{W}$, \qty{0.72}{W/kg}। छछूँदर:
$3.4\times 0.005^{0.75} = \qty{0.064}{W}$, \qty{13}{W/kg}। प्रति दिन:
$3\times 0.064\times 86400 = \qty{16.6}{kJ}$, यानी \qty{3.3}{g} भोजन,
अपने ही द्रव्यमान का दो-तिहाई।
\end{solution}

\begin{solution}{exo:b1:organism-environment:7}
$\qty{80}{W} = \qty{4.8}{kJ/min}$, यानी प्रति मिनट $\mathrm{O_2}$ की
\qty{0.24}{L}; प्रति दिन $80\times 86400 = \qty{6.9}{MJ}$। \qty{2000}{kcal}
= प्रति दिन \qty{8.36}{MJ} = \qty{97}{W}।
\end{solution}

\begin{solution}{exo:b1:organism-environment:8}
चूहा: $100\times(0.02/70)^{0.75} = \qty{0.22}{m^2}$; हाथी:
$100\times(4000/70)^{0.75} = \qty{2080}{m^2}$। गोले: चूहा
$r = \qty{1.7}{cm}$, $S = \qty{0.0036}{m^2}$ (फेफड़ा शरीर के पृष्ठ का
साठ गुना); हाथी $r = \qty{0.98}{m}$, $S = \qty{12}{m^2}$ (फेफड़ा
एक सौ सत्तर गुना)। जंतु जितना बड़ा, उसका फेफड़ा अपने बाहरी पृष्ठ के
सापेक्ष उतना ही अधिक मुड़ा होना चाहिए।
\end{solution}

\begin{solution}{exo:b1:organism-environment:9}
प्रभावक तभी काम करते हैं जब संवेदी कोई विचलन दर्ज कर चुका हो, इसलिए किसी
विचलन का होना ज़रूरी है इससे पहले कि वह सुधारा जाए; और सुधार तब ज़रा आगे
निकल जाता है इससे पहले कि संवेदी उसकी सूचना दे। पट्टी की चौड़ाई संवेदी की
सुग्राहिता (सबसे छोटा पकड़ में आता विचलन), पकड़ और असर के बीच के विलंब,
तथा प्रभावकों की ताकत से तय होती है।
\end{solution}

\begin{solution}{exo:b1:organism-environment:10}
$V = \qty{2e-5}{m^3}$, $r = \qty{1.68}{cm}$, $S = 4\pi r^2 =
\qty{3.6e-3}{m^2}$। हानि: $5\times 3.6\times 10^{-3}\times 37 =
\qty{0.66}{W}$। आधारी उत्पादन: \qty{0.18}{W}, यानी चार गुना कम। चूहे को
अपना उत्पादन चढ़ाना (कँपकँपी, भोजन: आधारी का तीन-चार गुना), $h$ घटाना
(घोंसला, चिपककर बैठना) या अपना तापमान गिरने देना (सुषुप्ति) पड़ता है।
\end{solution}

\begin{solution}{exo:b1:organism-environment:11}
हानि $\propto M^{2/3}$, उत्पादन $\propto M^{3/4}$: इसलिए अनुपात
उत्पादन/हानि $\propto M^{1/12}$ द्रव्यमान के साथ बढ़ता है, यानी छोटे
जंतुओं का प्रति इकाई पृष्ठ उत्पादन सबसे कम और बड़ों का सबसे अधिक होता है।
$M$ द्रव्यमान का गोला: $S = 4\pi(3M/4000\pi)^{2/3} = \num{0.0483}\,M^{2/3}$ (SI);
हानि $= 5\times 20\times 0.0483\,M^{2/3} = 4.83\,M^{2/3}$; और
$3.4\,M^{3/4}$ के साथ बराबरी से $M^{1/12} = 1.42$, $M \approx \qty{67}{kg}$।
इससे नीचे विश्राम करता जंतु बिना कुचालन के ठंडा रहता है; इससे ऊपर वह गरम
हो जाता है और उसे ऊष्मा बहानी पड़ती है।
\end{solution}

\begin{solution}{exo:b1:organism-environment:12}
अधिक सरल नहीं, बल्कि अलग ढंग से संगठित। पौधे के विनिमय बाहरी हैं (पत्ती
और जड़ के पृष्ठ माध्यम में फैले हुए) और उसकी कोशिकाएँ ऐसे माध्यम के सामने
खुली रहती हैं जिसका वह नियमन नहीं करता; उसका ``नियमन'' वृद्धि और रूप है,
जो पृष्ठों को ही समायोजित करता है। स्तनधारी अपने विनिमय किसी नियमित तरल
के पीछे भीतर कर लेता है ताकि उसकी कोशिकाएँ बाहर देखें ही नहीं। \omterm{def:b1:organism-environment:homeostasis}{समस्थापन}
जंतु के आंतरिक परिवेश का गुण है, जटिलता की माप नहीं: पौधे की कोशिकाएँ
अपनी ही सामग्री का नियमन करती हैं, उसके रंध्र उसकी जल-हानि का नियमन करते
हैं, और उसका शरीर लगातार बढ़ता हुआ एक विनिमय-पृष्ठ है --- यह एक रणनीति
है, कोई अभाव नहीं।
\end{solution}

\begin{solution}{pb:b1:organism-environment:1}
\textbf{1.} $V = M/\rho$: चूहा \qty{2e-5}{m^3} (\qty{20}{cm^3}),
हाथी \qty{4}{m^3}।
\textbf{2.} $r = (3V/4\pi)^{1/3}$: चूहा \qty{1.68}{cm}, हाथी
\qty{0.985}{m}।
\textbf{3.} $S = 4\pi r^2$: चूहा \qty{3.55e-3}{m^2} (\qty{35.5}{cm^2}),
हाथी \qty{12.2}{m^2}।
\textbf{4.} $S/V = 3/r$: चूहा \qty{178}{m^{-1}}, हाथी
\qty{3.05}{m^{-1}}; अनुपात $58$।
\textbf{5.} $S/V \propto 1/r \propto M^{-1/3}$, इसलिए अनुपात
$(M_e/M_m)^{1/3} = (2\times 10^5)^{1/3} = 58.5$ है।
\textbf{6.} दोनों के लिए बड़ा (किसी गोले का पृष्ठ अपने आयतन के लिए सबसे
छोटा होता है); चूहा, अपनी पूँछ, कानों और पतले अंगों के साथ, अधिक हटता
है।
\textbf{7.} $P = 10\times S\times 20$: चूहा \qty{0.71}{W}, हाथी
\qty{2440}{W}।
\textbf{8.} चूहा $3.4\times 0.02^{0.75} = \qty{0.18}{W}$; हाथी
$3.4\times 4000^{0.75} = \qty{1710}{W}$।
\textbf{9.} विश्राम पर \qty{1710}{W} उत्पादन के सामने \qty{2440}{W}
हानि: नंगी त्वचा मोटे तौर पर बहुत है, और हाथी को सक्रिय होने पर या गरम
हवा में ठंडा होने से अधिक गरम हो जाने का खतरा रहता है।
\textbf{10.} \qty{0.18}{W} के सामने \qty{0.71}{W} हानि: चार गुना अधिक।
रोएँ को $h$ घटाकर लगभग \qty{2.5}{W.m^{-2}.K^{-1}} पर लाना पड़ता है।
\textbf{11.} उसका $S/V$ किसी भी थलचर स्तनधारी से छोटा है; कान उसके पृष्ठ
में एक-तिहाई जोड़ते हैं और, पतले तथा रक्त से भरे होने के कारण, ऐसे
विकिरक हैं जो तब ऊष्मा फेंकते हैं जब शरीर \qty{12}{m^2} की हानि-क्षमता से
अधिक बना रहा हो।
\textbf{12.} $S/V \propto M^{-1/3}$ तब बिना सीमा के बढ़ता है जब
$M$ गिरे, जबकि $B/M \propto M^{-1/4}$ अधिक धीरे बढ़ता है; कुछ ग्राम से नीचे न
कोई रोएँ और न कोई भोजन-दर उस हानि को ढँक सकती है। रात में तापमान गिरने
देना (सुषुप्ति) बिना भोजन वाले घंटों के लिए $T_b - T_a$ पद ही हटा देता है।
\textbf{13.} $2B\times\qty{86400}{s}$: चूहा \qty{31}{kJ}, हाथी
\qty{295}{MJ}।
\textbf{14.} चूहा $31/18 = \qty{1.7}{g}$; हाथी
$295\,000/5 = \qty{59}{kg}$।
\textbf{15.} चूहा अपने द्रव्यमान का $1.7/20 = 8.6\%$; हाथी $59/4000 =
1.5\%$: चूहा प्रति ग्राम छह गुना अधिक खाता है।
\textbf{16.} आँत की लंबाइयों का अनुपात $35/0.4 = 88$; त्रिज्याओं का
अनुपात $0.985/0.0168
= 59$। आँत शरीर की रैखिक माप से तेज़ बढ़ती है (यहाँ
मोटे तौर पर $M^{0.36}$ की तरह), जिससे हाथी का अवशोषक पृष्ठ उसके आयतन के
सापेक्ष चढ़ जाता है।
\textbf{17.} $f_0 = 600\times 0.02^{1/4} = \qty{226}{min^{-1}}$;
$\tau_0 = 2/0.02^{1/4} = \qty{5.3}{yr}$।
\textbf{18.} $f = 226\times 4000^{-1/4} = \qty{28}{min^{-1}}$; $\tau =
5.3\times 4000^{1/4} = \qty{42}{yr}$।
\textbf{19.} चूहा $600\times 60\times 24\times 365\times 2 =
\num{6.3e8}$; हाथी $28\times 60\times 24\times 365\times 42 =
\num{6.2e8}$।
\textbf{20.} प्रति जीवन धड़कनें $= f\tau \propto M^{-1/4}M^{1/4} = M^0$:
द्रव्यमान से स्वतंत्र, लगभग एक अरब।
\textbf{21.} $f = 226\times 70^{-1/4} = \qty{78}{min^{-1}}$ (70 के
पास); $\tau = 5.3\times 70^{1/4} = \qty{15}{yr}$, जो 80 से बहुत नीचे है:
मानव जीवनकाल ही विषम है, स्तनधारी प्रवृत्ति का तीन से पाँच गुना, और उसके
कारण (मस्तिष्क, सामाजिकता, चिकित्सा) मापनी-नियम के बाहर हैं।
\textbf{22.} $160\times 4000^{-1/4}/0.02^{-1/4} = 160\times 0.0473 =
\qty{7.6}{min^{-1}}$; प्रति साँस धड़कनें $600/160 = 28/7.6 = 3.75$
दोनों के लिए, क्योंकि दोनों दरें एक ही घातांक साझा करती हैं।
\textbf{23.} $B/M = 3.4\,M^{3/4}/M = 3.4\,M^{-1/4}$; घातांक $-1/4$।
\textbf{24.} चूहा $0.18/0.02 = \qty{9.0}{W/kg}$; हाथी $1710/4000 =
\qty{0.43}{W/kg}$।
\textbf{25.} अनुपात $9.0/0.43 = 21$, जो $(M_e/M_m)^{1/4} =
(2\times 10^5)^{1/4} = 21.1$ के बराबर है: चूहे का हर ग्राम हाथी के हर
ग्राम से इक्कीस गुना तेज़ जलता है --- यानी द्रव्यमान-विशिष्ट उपापचयी
अनुपात द्रव्यमान-अनुपात का चौथा मूल है।
\end{solution}
